\documentclass[10pt,a4paper]{article}
\usepackage{amsmath,amssymb,amsthm,graphicx}
\usepackage{hyperref}
\usepackage{mathtools}
\usepackage{enumitem}
\DeclareMathOperator{\bias}{bias}
\DeclareMathOperator{\Cov}{Cov}
\DeclareMathOperator{\dsum}{dsum}
\title{Binary codes closer to the Gilbert-Varshamov bound admit near-linear unique and list decoding}
\author{Guy Blanc, Dean Doron}
\date{doi:10.4230/LIPIcs.CCC.2022.10; CC BY 4.0}
\begin{document}
\maketitle
\noindent\textit{Editorial note.} Complete selected author published code proof faithfully transcribed from CC BY4 formal version, including all local bearing support and AppendixA. Original mathematical expressions, hypotheses and numbered labels preserved; no mathematical repairs or new assumptions supplied. Only display line breaks and ordinary mathematical typography adapt presentation. Reused standard external decoder/base-code results are quoted at the exact original interface. All production and validation solely onYupeng viaSSH.\par\medskip
\noindent\textit{Allowed standard prerequisites (editorial).}\begin{itemize}
\item Bilu-Linial inverse mixing lemma, Hoeffding inequality and explicit ordinary expanders. Exact original cited matrixLemma16 and concentrationFact19 retained; ordinary expander existence is cited as by authors, not recursively reproved
\item Guruswami-Indyk base codes and Jeronimo-Srivastava-Tulsiani decoding framework. Original exact Theorems30/35 retained as cited standard inputs; full source-local compatibilityAppendixA and all new sampling/parameter arguments included
\end{itemize}\medskip
\section*{2 Preliminaries}
For integers $a,b$, we use $[a,b]$ to denote the set $\{a,\ldots,b\}$ and $[n]$ as a shorthand for $[1,n]$. Given a set $S\subseteq[n]$, when the ground set $[n]$ is clear from context, we denote $\rho(S)=|S|/n$, and its $\pm1$ variant as $\bias(S)=1-2\rho(S)$. For $z\in\mathbb F_2^n$, we similarly denote $\bias(z)$ as the bias of its characteristic set, i.e., $\mathbb E_{i\in[n]}[(-1)^{z_i}]$. We use boldface letters to denote random variables, except for $\mathbf1\in\mathbb R^n$, which we use for the all-ones vector. Also, when bounding running time, by writing $g(n)=\exp(f(n))$ we mean that $g(n)\leq2^{c\cdot f(n)}$ for some universal constant $c>0$.
\paragraph{Definition5 (discretizable distribution).} For $M\in\mathbb N$, We say that a distribution $\mathbf W$ is $M$-discretizable if it satisfies either of the following two equivalent properties.
\begin{enumerate}
\item For any $x$ in the support of $\mathbf W$, $\Pr_{\mathbf x\sim\mathbf W}[\mathbf x=x]=i/M$ for some $i\in\mathbb N$.
\item Let $U_M$ be the uniform distribution over $[M]$. Then, there is some function $f$ mapping $[M]$ to the support of $\mathbf W$ for which $f(i)$, where $i\sim U_M$, has the same distribution as a sample from $\mathbf W$.
\end{enumerate}
We say that $\mathbf W$ is computable in (deterministic or probabilistic) time $T$ if $f$ above is computable in time $T$. In particular, $\mathbf W$ is explicit if it is computable in deterministic time $\operatorname{poly}(M)$, and fully explicit if it is computable in deterministic time $\operatorname{poly}(\log M)$.
\paragraph{Definition6 (homogeneous distribution).} For $n,t\in\mathbb N$, we say that a distribution $\mathbf W$ over $[n]^t$ is homogeneous if its restriction to any coordinate is uniform over $[n]$, i.e., if for any $i\in[t]$ and $a\in[n]$ it holds that $\Pr_{\mathbf w\sim\mathbf W}[\mathbf w_i=a]=1/n$.
For any domain $\mathcal X$ and two distributions $\mathbf D,\mathbf D'$ over $\mathcal X$ we define the total variation distance of $\mathbf D$ and $\mathbf D'$ in terms of the optimal test distinguishing the distributions, i.e.,
\[d_{\mathrm{TV}}(\mathbf D,\mathbf D')\coloneqq\sup_{T:\mathcal X\to[0,1]}\{\mathbb E_{\mathbf x\sim\mathbf D}[T(\mathbf x)]-\mathbb E_{\mathbf x\sim\mathbf D'}[T(\mathbf x')]\}.\]
For a matrix $A\in\mathbb R^{n\times n}$, we denote by $\|A\|_{\mathrm{op}}$ its operator norm $\|A\|_{\mathrm{op}}=\max_{x\neq0}\|Ax\|_2/\|x\|_2$, which is also the maximum of $x^{\dagger}Ay$ over all norm-1 vectors $x,y\in\mathbb R^n$.
\paragraph{Error correcting codes.}
A binary error correcting code of message length $k$ and blocklength $n$ is a mapping $C:\mathbb F_2^k\to\mathbb F_2^n$, which we will often identify with its image $\operatorname{Im}(C)\subseteq\mathbb F_2^n$. The rate of $C$ is $k/n$, and its relative distance is $\Delta(C)=\frac1n\min_{z\neq z'}\Delta(z,z')$ for $z,z'\in C$, and $\Delta(z,z')=|\{i\in[n]:z_i\neq z'_i\}|$ being the Hamming distance. The Hamming ball of (relative) radius $\beta$ centered at $z$ is the set $B(z,\beta)=\{z'\in\mathbb F_2^n:\Delta(z,z')/n\leq\beta\}$.
We denote $\bias(C)$ as the maximal bias, in absolute value, of every nonzero $z\in C$. Thus, $\bias(C)\leq\varepsilon$ if the Hamming weight of any nonzero codeword is in $[(1-\varepsilon)/2,(1+\varepsilon)/2]$.
\paragraph{Definition7 (balanced codes).} A linear binary error correcting code is $\varepsilon$-balanced if $\bias(C)\leq\varepsilon$. In particular, an $\varepsilon$-balanced code $C$ has distance at least $(1-\varepsilon)/2$.
\paragraph{Unique and list decoding.}
We say that $C$ is (combinatorially) $(\beta,L)$ list decodable if for every $z\in\mathbb F_2^n$, $|C\cap B(z,\beta)|\leq L$. The Johnson bound tells us that any $\varepsilon$-balanced code is $(1/2-\sqrt\varepsilon,L)$ list decodable for $L=O(1/\varepsilon)$. The algorithmic list decoding problem aims at finding the list $L_{C,\beta}(z)\coloneqq C\cap B(z,\beta)$. When $\beta\leq(1-\varepsilon)/4$ and $C$ is $\varepsilon$-balanced, we know that $L_{C,\beta}$ always contains at most one codeword, which corresponds to the unique decoding problem.
\subsection*{2.1 Parity samplers and direct sum codes}
We will be interested in constructing a binary linear code $C\subseteq\mathbb F_2^{\bar n}$ with small bias by amplifying the (moderate) bias of some base code $C_0:\mathbb F_2^k\to\mathbb F_2^n$. One natural way to do so is by XORing $t$-tuples of $C_0$ according to some distribution $\mathbf W\sim[n]^t$.
\paragraph{Definition8 (direct sum codes).} For $t,n,\bar n\in\mathbb N$, let $\mathbf W\sim[n]^t$ be an $\bar n$-discretizable distribution equipped with a corresponding mapping function $f_{\mathbf W}:[\bar n]\to[n]^t$. For $z\in\mathbb F_2^n$, we let $\dsum_{\mathbf W}(z)\in\mathbb F_2^{\bar n}$ be such that
\[\dsum_{\mathbf W}(z)[\ell]=\sum_{i=1}^t z[f_{\mathbf W}(\ell)_i]\]
where the addition is taken over $\mathbb F_2$. Given a code $C_0\subseteq\mathbb F_2^n$, the direct sum lift of $C_0$ according to $\mathbf W$ is the code $\dsum_{\mathbf W}(C_0)=\{\dsum_{\mathbf W}(z):z\in C_0\}\subseteq\mathbb F_2^{\bar n}$.
\paragraph{Definition9 (parity sampler).} For $t,n\in\mathbb N$, and $0\leq\varepsilon<\varepsilon_0\leq1$, we say that $\mathbf W\sim[n]^t$ is an $(\varepsilon_0,\varepsilon)$ parity sampler if for every $z\in\mathbb F_2^n$ with $|\bias(z)|\leq\varepsilon_0$ it holds that $|\bias(\dsum_{\mathbf W}(z))|\leq\varepsilon$.
Clearly, if $C_0:\mathbb F_2^k\to\mathbb F_2^n$ is $\varepsilon_0$-balanced and $\mathbf W\sim[n]^t$ is an $\bar n$-discretizable $(\varepsilon_0,\varepsilon)$ parity sampler, the lifted code $C=\dsum_{\mathbf W}(C_0)$ is $\varepsilon$-balanced with rate $k/\bar n$.
\section*{3 Expanding 3-Uniform Hypergraphs}
Our construction uses a family of expanding $d$-regular3-uniform hypergraphs.
\paragraph{Definition10 ($d$-regular3-uniform hypergraph).} A3-uniform hypergraph consists of a set of vertices, $V$, and hyperedges $E\subseteq V^3$. The hypergraph $H=(V,E)$ is $d$-regular if, for each $v\in V$ and $j\in[3]$, the number of hyperedges $(w_1,w_2,w_3)\in E$ for which $v=w_j$ is $d$.
We will set $d$ carefully in our construction, but for now, it can be thought of as an arbitrary constant. For the remainder of this section, we will use ``hypergraph'' as shorthand for $d$-regular3-uniform hypergraph. There are various notions of expansion for hypergraphs and they are not all equivalent. In this work, we consider two such notions.
\paragraph{Definition11 ($\lambda$-mixing hypergraph).} A $d$-regular hypergraph $H=(V,E)$ on $n$ vertices is $\lambda$-mixing if, for any $S_1,S_2,S_3\subseteq V$,
\[\left|\frac{E(S_1,S_2,S_3)}d-\frac{|S_1|\cdot|S_2|\cdot|S_3|}{n^2}\right|\leq\lambda\sqrt{|S_1|\cdot|S_3|},\tag{2}\]
where $E(S_1,S_2,S_3)$ is the number of hyperedges $(w_1,w_2,w_3)\in E$ where $w_j\in S_j$ for $j=1,2,3$.
Note that the right-hand side of the above definition does not depend on the size of $S_2$.\footnote{Some works consider the stronger requirement of $\sqrt{|S_{\sigma(1)}|\cdot|S_{\sigma(2)}|}$ instead of $\sqrt{|S_1|\cdot|S_3|}$, for $S_{\sigma(1)}$ and $S_{\sigma(2)}$ being the two smallest sets.} That parallels the definition below, in which we use $\|y\|_{\infty}$ instead of $\|y\|_2$.
\paragraph{Definition12 ($\lambda$-spectral hypergraph).} A $d$-regular hypergraph $H=(V,E)$ on $n$ vertices is $\lambda$-spectral if, for any $x,y,z\in\mathbb R^n$,
\begin{multline*}
\left|\frac1d\sum_{(i,j,k)\in E}x_iy_jz_k-\frac1{n^2}\sum_{i\in V}x_i\sum_{j\in V}y_j\sum_{k\in V}z_k\right|\\
\leq\lambda\|x\|_2\|y\|_{\infty}\|z\|_2.\tag{3}
\end{multline*}
We will only care about cases where $y\in\{\pm1\}^n$.\footnote{In fact, for any fixed choice of $x$ and $z$, the left-hand side of Equation(3) is linear in $y$. Therefore, it is maximized for some $y\in\{\pm1\}^n$ and so in general it is sufficient to consider only such $y$.} In those cases, we have $\|y\|_{\infty}=1$ and $\|y\|_2=\sqrt n$. It is thus tempting to replace the $\|y\|_{\infty}$ in the right-hand side of Equation(3) with $\|y\|_2/\sqrt n$, as $\ell_2$ norms are often easier to work with than $\ell_{\infty}$ norms. Unfortunately, no ``good'' $\lambda$-spectral hypergraphs would exist with that modification: Whenever $n\gg d^2$, it is straightforward to bound $\lambda=\Omega(\sqrt n/d)$. For our definition, as we shall soon see, it is possible to achieve $\lambda\approx1/\sqrt d$.
Similar to standard graphs, spectral expansion implies mixing.
\paragraph{Proposition13 (spectral implies mixing).} For any $\lambda>0$, if $H$ is a $\lambda$-spectral hypergraph, it is also a $\lambda$-mixing hypergraph.
\begin{proof}
For any $S_1,S_2,S_3\subseteq V$, let $x_i=1[i\in S_1]$, $y_i=1[i\in S_2]$, and $z_i=1[i\in S_3]$. Then, $E(S_1,S_2,S_3)=\sum_{(i,j,k)\in E}x_iy_jz_k$. Equation(2) follows directly from Equation(3).
\end{proof}
By applying the converse to the expander mixing lemma for ordinary graphs[7], we can show that for symmetric hypergraphs, mixing implies spectral expansion with only a minor quantitative gap.
\paragraph{Definition14 (symmetric3-uniform hypergraph).} We say a3-uniform hypergraph $H=(V,E)$ is symmetric if, for any edge $e=(v_1,v_2,v_3)\in E$, the edge $(v_3,v_2,v_1)$ is also in $E$.

\paragraph{Theorem31.} There exists an efficient randomized algorithm such that for every$k$ and any$\varepsilon>0$, outputs with probability$1-2^{-\Omega(k)}$ an$\varepsilon$-balanced linear code$C:\mathbb F_2^k\to\mathbb F_2^{\bar n}$ of rate
\[\frac k{\bar n}=2^{-c_r\log\log(1/\varepsilon)\sqrt{\log(1/\varepsilon)}}\cdot\varepsilon^2\]
for some universal constant$c_r$, that is encodable in deterministic time$\exp(\exp(\sqrt{\log(1/\varepsilon)}))\cdot k$.\footnote{Following the footnote on the base code, one can get a randomized encoding in time$\operatorname{poly}(1/\varepsilon)\cdot k$ by using a randomized encoding of the based code.} More precisely, there exists a randomized preprocessing step that runs in time$\exp(\sqrt{\log(1/\varepsilon)})\cdot k$ and succeeds with probability$1-2^{-\Omega(k)}$. Once it succeeds, it fixes a deterministic mapping$C$ such that for every$x\in\mathbb F_2^k$, $C(x)$ can be computed in deterministic in the above$O_{\varepsilon}(\bar n)$ time.
If, moreover, for a large enough constant$d$ we are given an explicit family of$\lambda$-spectral$d$-regular3-uniform hypergraphs satisfying$\lambda=\operatorname{poly}(\log d)/\sqrt d$, there is no need for a preprocessing step, and$C$ is explicit.
\paragraph{Theorem36.} Given$k\in\mathbb N$ and$\varepsilon>0$, let$C:\mathbb F_2^k\to\mathbb F_2^{\bar n}$ be the$\varepsilon$-balanced, linear-time encodable code guaranteed to us by Theorem31 with high probability. Then,$C$ also admits the following decoding capabilities.
\begin{enumerate}
\item $C$ is list decodable up to radius$1/2-\beta$ for$\beta=2^{-\frac12\sqrt{\log(1/\varepsilon)}}$ by a randomized algorithm that runs in time$c_1(\varepsilon)\cdot\widetilde O(k)$ and succeeds with probability$1-2^{-\Omega(k)}$.
\item $C$ is uniquely decodable to within distance$(1-\varepsilon)/4$ by a randomized algorithm that runs in time$c_2(\varepsilon)\cdot\widetilde O(k)$ and succeeds with probability$1-2^{-\Omega(k)}$.
\end{enumerate}
Above,$c_1(\varepsilon)=\exp(\exp(\exp(\sqrt{\log(1/\varepsilon)})))$ and$c_2(\varepsilon)=\exp(\exp(\sqrt{\log(1/\varepsilon)}))$.

\paragraph{Proposition15 (mixing implies spectral).} There exists a universal constant $c_{\mathrm{spec}}>0$ for which the following holds. Let $H=(V=[n],E)$ be any $d$-regular hypergraph, and $\lambda\leq1/2$. If $H$ is $\lambda$-mixing and symmetric, then $H$ is a $\lambda'$-spectral for $\lambda'=c_{\mathrm{spec}}\cdot\lambda\log(1/\lambda)$.
The proof of Proposition15 is a simple application of a similar result for graphs.
\paragraph{Lemma16 (Lemma3.3 of[7]).} There exists a universal constant $c_{\mathrm{BL}}$ for which the following hods. For any $n\times n$ real symmetric matrix $A$ and $\lambda\leq1/2$, suppose each row has an $\ell_1$ norm of at most1 and for any two vectors $u,v\in\{0,1\}^n$,\footnote{Note that Bilu and Linial's Lemma statement only has the weaker requirement that this hold for orthogonal $u$ and $v$. As a result, they have an additional condition that the diagonal entries of $A$ not be too large, which is not needed for our version of the Lemma.}
\[|u^{\dagger}Av|\leq\lambda\|u\|\|v\|.\tag{4}\]
Then, the spectral radius of $A$ is at most $c_{\mathrm{BL}}\cdot\lambda\log(1/\lambda)$.
If $G$ is the (normalized) transition matrix of a $d$-regular graph, and $A=G-\frac1nJ$ for $J$ being the all-ones matrix, then Lemma16 shows a converse to the expander mixing lemma: Any graph with good mixing is also a good spectral expander. We show that their result can be lifted to3-uniform hypergraphs.
\begin{proof}[Proof of Proposition15]
Fix any $y\in\mathbb R^n$ and let $A^{(y)}\in\mathbb R^{n\times n}$ be defined as
\[A^{(y)}_{i,k}\coloneqq\frac1{2d}\left(\sum_{j\in[n]}1[(i,j,k)\in E]\cdot y_j\right)-\frac1{2n^2}\sum_{j\in[n]}y_j.\]
Then, for any $x,z\in\mathbb R^n$,
\[2(z^{\dagger}A^{(y)}x)=\frac1d\sum_{(i,j,k)\in E}x_iy_jz_k-\frac1{n^2}\sum_{i\in V}x_i\sum_{i\in V}y_i\sum_{i\in V}z_j.\]
Therefore, in order to prove that $H$ is an $\lambda'$-spectral expander, it is sufficient to show that for all $y\in\mathbb R^n$, the operator norm of $A^{(y)}$ is at most $\|y\|_{\infty}\cdot\lambda'/2$. We observe that:
\begin{enumerate}
\item For any fixed $x,z\in\mathbb R^n$, the quantity $z^{\dagger}A^{(y)}x$ is a linear function of $y$. Thus, we can consider $y$-s with $\|y\|_{\infty}=1$ without loss of generality. Furthermore, the $y$ maximizing $z^{\dagger}A^{(y)}x$ among those with $\|y\|_{\infty}=1$ will be in $\{\pm1\}^n$. Therefore we assume $y\in\{\pm1\}^n$ also without loss of generality.
\item Since $H$ is symmetric, the operator norm and spectral radius of $A^{(y)}$ are equal, so we instead bound the spectral radius.
\end{enumerate}
We will apply Lemma16 to each $A^{(y)}$. Note that:
\begin{itemize}
\item $A^{(y)}$ is symmetric, which follows immediately from the fact that $H$ is symmetric.
\item The $\ell_1$ norm of each row of $A^{(y)}$ is bounded by1:
\begin{multline*}
\sum_{k\in[n]}\left|\frac1{2d}\left(\sum_{j\in[n]}1[(i,j,k)\in E]\cdot y_j\right)+\frac1{2n^2}\sum_{j\in[n]}y_j\right|\\
\leq\frac1{2d}\cdot d+\frac1{2n^2}\cdot n\leq1.
\end{multline*}
\item Finally, fix some $u,v\in\{0,1\}^n$, and define the sets
\[S_1=\{i\in V:u_i=1\},\quad S_3=\{i\in V:v_i=1\},\quad S_2^+=\{i\in V:y_i=1\},\quad S_2^-=\{i\in V:y_i=-1\}.\]
Then,
\begin{align*}
|u^{\dagger}A^{(y)}v|&=\frac12\left|\frac1d(E(S_1,S_2^+,S_3)-E(S_1,S_2^-,S_3))\right.\\
&\hspace{3em}\left.-\frac{|S_1|\cdot|S_3|\cdot(|S_2^+|-|S_2^-|)}{n^2}\right|\\
&\leq\frac12\left|\frac1dE(S_1,S_2^+,S_3)-\frac{|S_1|\cdot|S_3|\cdot|S_2^+|}{n^2}\right|\\
&\quad+\frac12\left|\frac1dE(S_1,S_2^-,S_3)-\frac{|S_1|\cdot|S_3|\cdot|S_2^-|}{n^2}\right|\\
&\leq\lambda\sqrt{|S_1|\cdot|S_3|},
\end{align*}
where the last inequality follows from fact that $H$ is $\lambda$-mixing (applied to both expressions).
\end{itemize}
Thus, $|u^{\dagger}A^{(y)}v|\leq\lambda\|u\|\|v\|$ and we can apply Lemma16 to obtain $\|A^{(y)}\|_{\mathrm{op}}=c_{\mathrm{BL}}\cdot\lambda\log(1/\lambda)$, implying that $H$ is $(\lambda'=2c_{\mathrm{BL}}\cdot\lambda\log(1/\lambda))$-spectral.
\end{proof}

\subsection*{3.1 Random hypergraphs mix well}
In this section, we'll show that given an expanding graph $G$, its random hypergraph completion is, with high probability, a good expander. Throughout, we say that an undirected regular graph $G$ is a $\lambda$-expander if the second largest eigenvalue of its normalized adjacency matrix, in magnitude, is at most $\lambda$.
\paragraph{Definition17 (random hypergraph completion of a graph).} Let $G=(V=[n],E_G)$ be a $d$-regular graph. To sample a random hypergraph completion, $\mathbf H$ of $G$, we choose a uniformly random ordering of $G$'s edges, $\{(u_1,v_1),\ldots,(u_{nd},v_{nd})\}=E_G$. Then, we set $\mathbf H=(V,E_{\mathbf H})$, where
\[E_{\mathbf H}\coloneqq\{(u_i,\lceil i/d\rceil,v_i)\mid i\in[nd]\}.\]
Equivalently, for each $(u,v)\in E_G$, we choose an independent and uniform $\mathbf w\in V$ and add the hyperedge $(u,\mathbf w,v)$ to $E_{\mathbf H}$, conditioned on the resulting hypergraph $\mathbf H$ being $d$-regular.
Next, we prove that the random hypergraph completion mixes well with high probability.
\paragraph{Lemma18.} Let $G=(V=[n],E)$ be a $d$-regular $\lambda$-expander and $\mathbf H$ be a random hypergraph completion of $G$. With probability at least $1-2^{-n}$, $\mathbf H$ is a $\lambda'$-mixing hypergraph for $\lambda'=2\lambda+2/\sqrt d$.
In order to prove Lemma18, we'll use Hoeffding's inequality for sampling without replacements.
\paragraph{Fact19 (Hoeffding's inequality,[18]).} For any integers $a,k\leq m$, suppose there are $m$ items, of which $k$ of them are marked. Let $\mathbf x$ be a random variable indicating the number of marked items when $a$ of the $m$ items are sampled uniformly and independently without replacement. Then, for any $t\geq0$,
\[\Pr[|\mathbf x-(k/m)\cdot a|\geq t]\leq2\exp(-2t^2/a).\]
\begin{proof}[Proof of Lemma18]
Fix arbitrary $S_1,S_2,S_3\subseteq V$. We will show that with probability at least $1-2^{-4n}$ it holds that
\[\left|\frac{E_{\mathbf H}(S_1,S_2,S_3)}d-\frac{|S_1|\cdot|S_2|\cdot|S_3|}{n^2}\right|\leq\left(\frac2{\sqrt d}+2\lambda\right)\sqrt{|S_1|\cdot|S_3|}\tag{5}\]
where $E_{\mathbf H}(S_1,S_2,S_3)$ is the number of edges $(v_1,v_2,v_3)$ in $\mathbf H$ where $v_j\in S_j$ for each $j\in[3]$. The desired result then follows from a union bound over the $(2^n)^3=2^{3n}$ choices for $S_1,S_2,S_3$.
Let $E_G(S_1,S_3)$ be the number of edges, $(u,v)$, of $G$, such that $u\in S_1$ and $v\in S_3$. By the expander mixing lemma applied to $G$, we have that
\[\left|\frac{E_G(S_1,S_3)}d-\frac{|S_1||S_3|}n\right|\leq\lambda\sqrt{|S_1||S_3|}.\]
Let us define $\mu\coloneqq|S_1||S_3|/n$ and $\Delta\coloneqq\lambda\sqrt{|S_1||S_3|}$. We consider two cases.
\begin{enumerate}
\item In the first case, $\mu\leq\Delta$. Here, we use the simple bound
\[0\leq\frac{E_{\mathbf H}(S_1,S_2,S_3)}d\leq\frac{E_G(S_1,S_3)}d\leq2\Delta\]
that holds with probability1. This, along with the fact $|S_1|\cdot|S_2|\cdot|S_3|/n^2\leq\mu\leq\Delta$ implies that Equation(5) always holds.
\item In the second case, $\mu>\Delta$. Here, we will apply Fact19. $G$ has a total of $nd$ edges, of which $d\cdot|S_2|$ are matched to a vertex in $S_2$. $E_{\mathbf H}(S_1,S_2,S_3)$ samples $E_G(S_1,S_3)$ of those $nd$ edges (without replacement) and counts how many were among the $d\cdot|S_2|$ assigned to $S_2$. Hence, by Hoeffding's inequality, we have for any $t\geq0$,
\[\Pr\left[\left|E_{\mathbf H}(S_1,S_2,S_3)-\frac{d\cdot|S_2|}{nd}\cdot E_G(S_1,S_3)\right|\geq t\right]\leq2\exp\left(\frac{-2t^2}{E_G(S_1,S_3)}\right).\]
Then, setting $t=2\sqrt{d|S_1|\cdot|S_3|}$ and using the fact that $E_G(S_1,S_3)\leq d(\mu+\Delta)\leq2d\mu$,
\begin{multline*}
\Pr\left[\left|E_{\mathbf H}(S_1,S_2,S_3)-\frac{|S_2|}n\cdot E_G(S_1,S_3)\right|\geq2\sqrt{d|S_1|\cdot|S_3|}\right]\\
\leq2\exp\left(\frac{-8d|S_1|\cdot|S_3|}{2d\cdot(|S_1||S_3|/n)}\right)=2\exp(-4.5n)\leq2^{-4n}.
\end{multline*}
As the above shows, $E_{\mathbf H}(S_1,S_2,S_3)$ is within $\pm t$ of its expectation with probability at least $2^{-4n}$. When that occurs, we have that
\begin{align*}
&\left|\frac{E_{\mathbf H}(S_1,S_2,S_3)}d-\frac{|S_1|\cdot|S_2|\cdot|S_3|}{n^2}\right|\\
&=\frac1d\left|E_{\mathbf H}(S_1,S_2,S_3)-\frac{|S_2|d\mu}n\right|\\
&\leq\frac td+\frac{|S_2|}n\left|\frac{E_G(S_1,S_3)}d-\mu\right|\\
&\leq\frac td+\frac{|S_2|}n\cdot\Delta\qquad\text{(expander mixing lemma)}\\
&\leq2\sqrt{\frac{|S_1|\cdot|S_3|}d}+\lambda\sqrt{|S_1|\cdot|S_3|}.\qquad(|S_2|/n\leq1)
\end{align*}
\end{enumerate}
Hence, Equation(5) holds with probability at least $1-2^{-4n}$ in both cases, so we can union bound over the $2^{3n}$ choices for $S_1,S_2,S_3$.
\end{proof}
For our purposes, we will want the hypergraph to be symmetric. It is quite easy to ``symmetrize'' any hypergraph at only a modest cost to the degree.
\paragraph{Proposition20.} For any $n,d\in\mathbb N$, there is an algorithm running in time $O(nd)$ that takes as input any $d$-regular $\lambda$-mixing hypergraph $H=(V=[n],E_H)$ and outputs a $2d$-regular $\lambda$-mixing hypergraph $H'=(V=[n],E_{H'})$ over the same vertices.
\begin{proof}
For every edge $(u,v,w)\in E_H$ we include both $(u,v,w)$ and the reverse edge $(w,v,u)$ in $E_{H'}$.\footnote{Note that, we do this even if it results in duplicated hyperedges: If $(u,v,w)$ and $(w,v,u)$ are both in $E_H$, then there will be two copies of $(u,v,w)$ and two copies of $(w,v,u)$ in $E_{H'}$.} Clearly, this results in the degree of $H'$ being $2d$. Then, for any $S_1,S_2,S_3\subseteq V$,
\begin{align*}
&\left|\frac{E_{H'}(S_1,S_2,S_3)}{2d}-\frac{|S_1|\cdot|S_2|\cdot|S_3|}{n^2}\right|\\
&=\left|\frac{E_H(S_1,S_2,S_3)+E_H(S_3,S_2,S_1)}{2d}-\frac{|S_1|\cdot|S_2|\cdot|S_3|}{n^2}\right|\\
&\leq\frac12\left|\frac{E_H(S_1,S_2,S_3)}d-\frac{|S_1|\cdot|S_2|\cdot|S_3|}{n^2}\right|\\
&\quad+\frac12\left|\frac{E_H(S_3,S_2,S_1)}d-\frac{|S_1|\cdot|S_2|\cdot|S_3|}{n^2}\right|\\
&\leq\lambda\sqrt{|S_1|\cdot|S_3|}.
\end{align*}
Therefore, $H'$ is $\lambda$-mixing, as desired.
\end{proof}
We have explicit (and even fully explicit) constructions of Ramanujan graphs, i.e., $\lambda$-expanders for $\lambda\leq2\sqrt{d-1}/d$, albeit with some restrictions on $d$[26,28]. By manipulating Ramanujan graphs, Alon gave a construction of $d$-regular $\lambda$-expanders over $n$ vertices for any $d$ and $n$ while suffering only a tiny loss in $\lambda$ (see[1], and also[29,13] for weaker constructions). In particular, there exist explicit expanders with $\lambda=O(1/\sqrt d)$ for all $n$-s. We thus get the following corollary.
\paragraph{Corollary21.} There exists a probabilistic algorithm such that for any integer $n$ and even integer $6\leq d\leq n$, runs in time $\operatorname{poly}(n)$ and with probability at least $1-2^{-n}$ outputs a3-uniform symmetric $d$-regular hypergraph that is $\lambda=c_{\mathrm{rand}}/\sqrt d$-mixing, where $c_{\mathrm{rand}}\geq2$ is some universal constant.
We will refer to the above probabilistic construction as our preprocessing step.

\section*{4 From Hypergraphs to Parity Samplers}
Fix some $n,t,d\in\mathbb N$ and a $d$-regular3-uniform hypergraph $H=(V,E_H)$ over $n$ vertices. We will construct a parity sampler $\mathbf W_{H,t}\sim[n]^t$ from $H$ and show that when $H$ is a good spectral expander, $\mathbf W$ is a good parity sampler.
\paragraph{The construction.}
For each $v\in V$, there are $d$ edges $(v_1,v_2,v_3)\in E_H$ with $v_1=v$. For any $i\in[d]$, let $e_H(v,i)$ be the $i$th such edge (i.e., $e_H(v,i)_1=v$). Without loss of generality, we assume the vertices are each labeled with a unique integer between1 and$n$ (i.e $V=[n]$).
To sample $\mathbf w\sim\mathbf W_{H,t}$, independently sample a starting vertex $\mathbf v_0\in[n]$ and edge labels $\mathbf i_1,\ldots,\mathbf i_t\sim[d]$ uniformly. $\mathbf w$ will be a deterministic function of $\mathbf v_0$ and $\mathbf i_1,\ldots,\mathbf i_t$ computed as follows. For each $j=1,\ldots,t$,
\begin{enumerate}
\item Let $\mathbf e_j=e_H(\mathbf v_{j-1},\mathbf i_j)$.
\item Let $\mathbf v_j=(\mathbf e_j)_3$.
\item Let $\mathbf w_j=(\mathbf e_j)_2$.
\end{enumerate}
The sample is then $\mathbf w=(\mathbf w_1,\ldots,\mathbf w_t)$.
We present two simple claims about $\mathbf W_{H,t}$.
\paragraph{Claim23.} For any $t\in\mathbb N$ and a $d$-regular hypergraph $H$ over $n$ vertices, $\mathbf W_{H,t}$ is $(nd^t)$-discretizable.
\begin{proof}
The sample $\mathbf w\sim\mathbf W_{H,t}$ is a deterministic function of $\mathbf v_0$ and $\mathbf i_1,\ldots,\mathbf i_t$, and those variables are set to a uniform choice out of $nd^t$ possibilities. The claim then follows from Definition5.
\end{proof}
\paragraph{Claim24.} For any $t\in\mathbb N$, $d$-regular hypergraph $H$, and $j\in[t]$, $\mathbf e_j$ is uniform over all $nd$ edges of $H$. As a result, $\mathbf W_{H,t}$ is homogeneous.
\begin{proof}
We will first prove that each $\mathbf e_j$ is uniform over $E_H$ by induction on$j$. $\mathbf e_1$ is uniform over the $nd$ edges in $E_H$ as it is sampled by independently choosing a starting vertex $\mathbf v_0\sim V$ uniformly and then uniformly choosing one of its $d$-neighbors. Furthermore, for any $j\in[t-1]$, if $\mathbf e_j$ is uniform, then $\mathbf v_j$ is also uniform. Hence, $\mathbf e_{j+1}$ is sampled by selecting a uniform vertex and then (independently) one of its $d$-neighbors, so $\mathbf e_{j+1}$ is also uniform over $E_H$.
Next, we show $\mathbf W_{H,t}$ is homogeneous. Fix any $a\in[n]$ and $j\in[t]$. As $\mathbf w_j=a$ if and only if $\mathbf e_j$ is one of the $d$-edges in $E_H$ whose second vertex is$a$ and $\mathbf e_j$ is uniform over the $nd$ edges in $E_H$, $\Pr[\mathbf w_j=a]=d/(nd)=1/n$.
\end{proof}
Finally, we show that whenever $H$ is a good expander, $\mathbf W_{H,t}$ is a good parity sampler.
\paragraph{Theorem25.} For any $n,d,t\in\mathbb N$, $\lambda,\varepsilon_0>0$, and $H$ a $\lambda$-spectral $d$-regular3-uniform hypergraph on$n$ vertices, $\mathbf W_{H,t}$ is an $(\varepsilon_0,\varepsilon\coloneqq(\varepsilon_0+\lambda)^t)$-parity sampler.
As an immediate corollary of Theorem25 and Proposition15:
\paragraph{Corollary26.} There exists an absolute constant $c_{\mathrm{spec}}>0$ for which the following holds. For any $n,d,t\in\mathbb N$, $\lambda,\varepsilon_0>0$, and $H$ a $\lambda$-mixing symmetric $d$-regular3-uniform hypergraph on$n$ vertices, $\mathbf W_{H,t}$ is an $(\varepsilon_0,\varepsilon\coloneqq(\varepsilon_0+c_{\mathrm{spec}}\cdot\lambda\log(1/\lambda))^t)$-parity sampler.
Throughout the remainder of this section, we will use the shorthand $\mathbf W\coloneqq\mathbf W_{H,t}$. For proving Theorem25, it will be convenient to consider $\sigma\in\{\pm1\}^n$ rather than $z\in\mathbb F_2^n$ (as in Definition9) using the mapping $\sigma_i=(-1)^{z_i}$. Equivalent to Definition9, $\mathbf W$ is an $(\varepsilon_0,\varepsilon)$ parity sampler if $|\bias_{\mathbf W}(\sigma)|\leq\varepsilon$ for all $\sigma\in\{\pm1\}^n$ satisfying $|\bias(\sigma)|=|\mathbb E_{i\sim[n]}[\sigma_i]|\leq\varepsilon_0$, where
\[\bias_{\mathbf W}(\sigma)\coloneqq\mathbb E_{\mathbf w\sim\mathbf W}\left[\prod_{j=1}^t\sigma_{\mathbf w_j}\right].\]
Similarly to[34], we express that bias algebraically.
\paragraph{Lemma27.} Let $A^{(\sigma)}\in\mathbb R^{n\times n}$ be defined as
\[A^{(\sigma)}_{i,k}\coloneqq\frac1d\sum_{(i',j',k')\in E_H}\sigma_{j'}\cdot1[i'=i\land k'=k].\]
Then, $\bias_{\mathbf W}(\sigma)=\frac1n\mathbf1^{\dagger}(A^{(\sigma)})^t\mathbf1$.
\begin{proof}
Let $\mathbf v_0,\ldots,\mathbf v_t$ and $\mathbf w_1,\ldots,\mathbf w_t$ be the random variables defined in the construction of $\mathbf W$. We claim that for each $j\in\{0,1,\ldots,t\}$ and$v\in[n]$, that
\[\mathbb E\left[1[\mathbf v_j=v]\prod_{k=1}^j\sigma_{\mathbf w_k}\right]=\frac1n\left(\mathbf1^{\dagger}(A^{(\sigma)})^j\right)_v\tag{6}\]
By induction on$j$. Clearly, for$j=0$, Equation(6) holds as both sides are equal to$1/n$ for any$v\in[n]$. For$j\geq1$,
\begin{align*}
&\mathbb E\left[1[\mathbf v_j=v]\prod_{k=1}^j\sigma_{\mathbf w_k}\right]\\
&=\sum_{v'\in[n]}\mathbb E\left[1[\mathbf v_j=v\land\mathbf v_{j-1}=v']\prod_{k=1}^j\sigma_{\mathbf w_k}\right]\\
&=\sum_{v'\in[n]}\mathbb E\left[1[\mathbf v_{j-1}=v']\prod_{k=1}^{j-1}\sigma_{\mathbf w_k}\right]\mathbb E[1[\mathbf v_j=v]\sigma_{\mathbf w_j}\mid\mathbf v_{j-1}=v']\\
&=\sum_{v'\in[n]}\frac1n\left(\mathbf1^{\dagger}(A^{(\sigma)})^{j-1}\right)_{v'}\mathbb E[1[\mathbf v_j=v]\sigma_{\mathbf w_j}\mid\mathbf v_{j-1}=v']\\
&=\frac1n\sum_{(a,b,c)\in E_H}\left(\mathbf1^{\dagger}(A^{(\sigma)})^{j-1}\right)_a\cdot\frac1d\cdot\sigma_b\cdot1[v=c]\\
&=\frac1n\left(\mathbf1^{\dagger}(A^{(\sigma)})^j\right)_v,
\end{align*}
where the third equality is the inductive hypothesis. Then,
\begin{align*}
\bias_{\mathbf W}(\sigma)&=\mathbb E_{\mathbf w\sim\mathbf W}\left[\prod_{j=1}^t\sigma_j\right]\\
&=\sum_{v\in[n]}\mathbb E\left[1[\mathbf v_j=v]\prod_{k=1}^j\sigma_{\mathbf w_k}\right]=\frac1n\mathbf1^{\dagger}(A^{(\sigma)})^t\mathbf1.
\end{align*}
\end{proof}
Next, we use the fact that$H$ is a$\lambda$-spectral expander to reason about$A^{(\sigma)}$.
\paragraph{Proposition28.} Let$J_n\in\mathbb R^{n\times n}$ be the matrix in which every element is$1/n$. For any$\sigma\in\{\pm1\}^n$,
\[\|A^{(\sigma)}-\bias(\sigma)J_n\|_{\mathrm{op}}\leq\lambda.\]
\begin{proof}
Fix any$x,z\in\mathbb R^n$. Then,
\begin{align*}
z^{\dagger}(A^{(\sigma)}-\bias(\sigma)J_n)x&=\frac1d\sum_{(i,j,k)\in E_H}x_k\sigma_jz_i-\frac{\bias(\sigma)}n\sum_{i\in[n]}x_i\sum_{j\in[n]}z_j\\
&=\frac1d\sum_{(i,j,k)\in E_H}x_i\sigma_jz_k-\frac1{n^2}\sum_{k\in[n]}\sigma_k\sum_{i\in[n]}x_i\sum_{j\in[n]}z_j\\
&\leq\lambda\|x\|_2\|z\|_2.\qquad\text{(Definition12)}
\end{align*}
\end{proof}
As an immediate consequence, we have:
\paragraph{Corollary29.} For any$\sigma\in\{\pm1\}^n$, $\|A^{(\sigma)}\|_{\mathrm{op}}\leq|\bias(\sigma)|+\lambda$.
\begin{proof}
As$\|J_n\|_{\mathrm{op}}=1$, the desired result follows from the reversed triangle inequality applied to the operator norm.
\end{proof}
Finally, we prove Theorem25.
\begin{proof}[Proof of Theorem25]
For any$\sigma\in\{\pm1\}^n$ satisfying$|\bias(\sigma)|\leq\varepsilon_0$, we have:
\begin{align*}
|\bias_{\mathbf W}(\sigma)|&=\frac{\mathbf1^{\dagger}}{\sqrt n}(A^{(\sigma)})^t\frac{\mathbf1}{\sqrt n}\qquad\text{(Lemma27)}\\
&\leq\|(A^{(\sigma)})^t\|_{\mathrm{op}}\leq\|A^{(\sigma)}\|_{\mathrm{op}}^t\\
&\leq(|\bias(\sigma)|+\lambda)^t\leq(\varepsilon_0+\lambda)^t.\qquad\text{(Corollary29)}
\end{align*}
\end{proof}

\section*{5 Codes Closer to the GV Bound}
\subsection*{5.1 The construction}
We use our parity sampler to amplify the distance of a given base code via direct sum lifting. More formally, given$k\in\mathbb N$ and$\varepsilon>0$, let$\varepsilon_0=\varepsilon_0(\varepsilon)$ soon to be determined, and let
\begin{itemize}
\item $C_0:\mathbb F_2^k\to\mathbb F_2^n$ be an$\varepsilon_0$-balanced code, and,
\item $\mathbf W\sim[n]^t$ be an$(\varepsilon_0,\varepsilon)$-parity sampler which is$M$-discretizable.
\end{itemize}
Denote$\bar n=|\operatorname{Supp}(\mathbf W)|\leq M$. Thus, the lifted code$C=\dsum_{\mathbf W}(C_0)\subseteq\mathbb F_2^{\bar n}$ is clearly$\varepsilon$-balanced. For the base code, we use the codes by Guruswami and Indyk, that admit$O_{\varepsilon_0}(n)$-time encoding and decoding.
\paragraph{Theorem30 ([15]).} For every integer$n$ and any$\varepsilon_0>0$ there exists an$\varepsilon_0$-balanced code$C_0\subseteq\mathbb F_2^n$ of rate$\Omega(\varepsilon_0^3)$. Furthermore,$C_0$ is encodable in time$\exp(\operatorname{poly}(1/\varepsilon_0))n$ and decodable from$1/4-\varepsilon_0$ fraction of errors in time$\exp(\operatorname{poly}(1/\varepsilon_0))n$.\footnote{One can get a better randomized encoding time of$\operatorname{poly}(1/\varepsilon_0)$.}
\paragraph{Setting parameters.} Given$\varepsilon>0$, we henceforth set:
\begin{enumerate}
\item The initial bias of$C_0$ to$\varepsilon_0=\frac12\cdot2^{-\sqrt{\log(1/\varepsilon)}}$.
\item The number of steps over$H$ to$t=\lceil\sqrt{\log(1/\varepsilon)}\rceil$.
\item The degree$d$ of$H$ to be the smallest even integer for which
\[c_{\mathrm{spec}}\cdot c_{\mathrm{rand}}\cdot\frac1{\sqrt d}\log\left(\frac{\sqrt d}{c_{\mathrm{rand}}}\right)\leq\varepsilon_0.\]
This is chosen so that Corollary21 gives, with high probability, a symmetric hypergraph$H$ such that, by Corollary26,$\mathbf W_{H,t}$ is a$(\varepsilon_0,(2\varepsilon_0)^t)$-parity sampler.
\end{enumerate}
Using the above parameters in Theorem25, together with the random hypergraph of Corollary21. we get the following (randomized) error correcting code$C$ [Theorem31, stated above].
\begin{proof}[Proof of Theorem31]
Given$k\in\mathbb N$ and$\varepsilon>0$, we set$\varepsilon_0,t$, and$d$ as above. Theorem30 gives us a code$C_0:\mathbb F_2^k\to\mathbb F_2^n$ which is$\varepsilon_0$-balanced, for$n=O(k/\varepsilon_0^3)$. By Corollary21 we can output a$\lambda=c_{\mathrm{rand}}/\sqrt d$-mixing$H$ with probability at least$1-2^{-n}$ in polynomial time. This is our preprocessing step. By Corollary26 and our choice of parameters, indeed
\[\left(\varepsilon_0+c_{\mathrm{spec}}\cdot\lambda\log\frac1\lambda\right)^t\leq(2\varepsilon_0)^t\leq\varepsilon,\]
so$\mathbf W=\mathbf W_{H,t}$ appropriately amplifies the distance. Moreover,$\mathbf W$ is$M$-discretizable for
\begin{align*}
M&=nd^t=O\left(\frac k{\varepsilon_0^3}\right)\cdot d^{\lceil\sqrt{\log(1/\varepsilon)}\rceil}\\
&=k\cdot2^{(3+\log d)\sqrt{\log(1/\varepsilon)}+\log d+O(1)}.
\end{align*}
We proceed to bounding$\log d$ in terms of$\varepsilon$. Recalling how we set$d$, we see that$\frac c{\sqrt d}\log d=2^{-\sqrt{\log(1/\varepsilon)}}$, where$c$ is some universal positive constant. From this we can infer that
\[d\leq c^2 2^{2\sqrt{\log(1/\varepsilon)}}\cdot\log^3\left(c\cdot2^{\sqrt{\log(1/\varepsilon)}}\right),\tag{7}\]
so$\log d=2\sqrt{\log(1/\varepsilon)}+O(\log\log(1/\varepsilon))$. (We are assuming that$\varepsilon$ is smaller than some small constant, which we can without loss of generality.) Hence,
\[M=\frac k{\varepsilon^2}\cdot2^{O(\sqrt{\log(1/\varepsilon)}\cdot\log\log(1/\varepsilon))}.\]
Since$\mathbf W$ is$M$-discretizable,$C=\dsum_{\mathbf W}(C_0)$ has block length$\bar n\leq M$, as required. Drawing$H$ at random, by Corollary21 takes
\[O(dn)=\widetilde O\left(\frac1{\varepsilon_0^3}\right)\cdot k=\exp(\sqrt{\log(1/\varepsilon)})\cdot k\]
time. Given the hypergraph$H$, the encoding amounts to computing$C_0(x)$ and taking parities of its coordinates according to$\mathbf W$. This takes$\exp(\operatorname{poly}(1/\varepsilon_0))\cdot k+M\cdot t$ time, following Theorem30.
For the ``moreover'' part, one can verify that we can tolerate a$\operatorname{polylog}(d)$ multiplicative factor in$\lambda$ with no substantial loss in parameters.
\end{proof}

\subsection*{5.2 $\tau$-sampling}
Toward establishing efficient decoding, we will need the notion of$\tau$-sampling.
\paragraph{Definition32 ($\tau$-sampling).} We say that a distribution$\mathbf W$ over$[n]^t$ is$\tau$-sampling if for any$i\in[t-1]$, $S\subseteq[n]$, and$X\subseteq[n]^i$,
\[\Cov_{\mathbf w\sim\mathbf W}[1[\mathbf w_{i+1}\in S],1[(\mathbf w_1,\ldots,\mathbf w_i)\in X]]\leq\tau.\]
To highlight the fact that it is indeed a (strong) sampling property, assume for simplicity that$\mathbf W$ is homogeneous. Fix any$i\in[t-1]$, $S\subseteq[n]$, and$X\subseteq[n]^i$. The property of$\tau$-sampling thus tells us we can use$\mathbf W$ to sample$S$ starting from any prefix. Namely, that
\[\left|\Pr_{\mathbf w\in\mathbf W}[\mathbf w_{i+1}\in S\mid(\mathbf w_1,\ldots,\mathbf w_i)\in X]-\rho(S)\right|\leq\frac\tau{\rho(X)}.\]
Note that the$t=2$ case corresponds to the setting of the expander mixing lemma.
A$\tau$-sampling distribution$\mathbf W$ satisfies the property of the ``Splittable Mixing Lemma'' of Jeronimo, Srivastava, and Tulsiani [20,Lemma4.6] for the family of ``$\pm1$ cut functions''. This fact is crucial us, and for completeness we establish this in AppendixA.
Given a$\tau$-sampling homogeneous$\mathbf W\sim[n]^t$, a standard hybrid argument shows the following.
\paragraph{Lemma33.} Let$\mathbf W\sim[n]^t$ be homogeneous and$\tau$-sampling, and let$z\in\mathbb F_2^n$ be arbitrary. Then,
\[\left|\mathbb E_{\mathbf w\in\mathbf W}[(-1)^{z_{\mathbf w_1}+\cdots+z_{\mathbf w_t}}]-\bias(z)^t\right|\leq2(t-1)\tau.\]
Thus, high-order mixing in particular implies that$\mathbf W$ is an$(\varepsilon_0,\varepsilon=\varepsilon_0^t+(t-1)\tau)$ parity sampler for all$\varepsilon_0$. However, for us, and also in[20], $t\cdot\tau$ is too large so we prove the parity sampling property separately.
We conclude our discussion about$\tau$-sampling by showing that our$\mathbf W$ is indeed$\tau$-sampling.
\paragraph{Lemma34.} For any$n,d,t\in\mathbb N$ and$\lambda>0$, if$H=(V=[n],E)$ is a$\lambda$-spectral$d$-regular3-uniform hypergraph, then$\mathbf W_{H,t}$ is$\tau=\lambda/4$-sampling.
\begin{proof}
As in Definition32, fix any$i\in[t-1]$, $S\subseteq[n]$ and$X\subseteq[n]^i$. Let$S_i(\mathbf w)$ be the indicator that$\mathbf w_{i+1}\in S$, and$X_i(\mathbf w)$ the indicator that$(\mathbf w_1,\ldots,\mathbf w_t)\in X$. Using the identity$\Cov[1-2a,1-2b]=4\Cov[a,b]$, it is sufficient to prove that
\[\Cov_{\mathbf w\sim\mathbf W_{H,t}}[(-1)^{S_i(\mathbf w)},(-1)^{X_i(\mathbf w)}]\leq\lambda.\]
Let$\mathbf v_0,\ldots,\mathbf v_t$, $\mathbf w_1,\ldots,\mathbf w_t$, and$\mathbf e_1,\ldots,\mathbf e_t$ be the random variables defined in the construction of$\mathbf W_{H,t}$. Furthermore, let$x,y,z\in\mathbb R^n$ be the vectors defined, for each$j\in[n]$, by
\[x_j\coloneqq\mathbb E[(-1)^{X_i(\mathbf w)}\mid\mathbf v_i=j],\qquad y_j\coloneqq(-1)^{1[j\in S]},\qquad z_j\coloneqq1/n.\]
Note that$\|z\|_2=1/\sqrt n$, $\|y\|_{\infty}=1$, and$\|x\|_2\leq\sqrt n$ (which follows from$\|x\|_{\infty}\leq1$). Our goal in this proof will be to show that the following equation holds:
\begin{multline*}
\Cov_{\mathbf w}[(-1)^{S_i(\mathbf w)},(-1)^{X_i(\mathbf w)}]\\
=\frac1d\sum_{(a,b,c)\in E}x_ay_bz_c-\frac1{n^2}\sum_{a\in[n]}x_a\sum_{a\in[n]}y_a\sum_{a\in[n]}z_a.\tag{8}
\end{multline*}
Once we do, the desired result follows from Definition12. In order to compute the covariance, we first expand:
\begin{align*}
&\mathbb E_{\mathbf w}[(-1)^{S_i(\mathbf w)}\cdot(-1)^{X_i(\mathbf w)}]\\
&=\frac1{nd}\sum_{e\in E}\mathbb E_{\mathbf w}[(-1)^{S_i(\mathbf w)}\cdot(-1)^{X_i(\mathbf w)}\mid\mathbf e_{i+1}=e]\qquad\text{(Claim24)}\\
&=\frac1{nd}\sum_{(a,b,c)\in E}(-1)^{1[b\in S]}\cdot\mathbb E_{\mathbf w}[(-1)^{X_i(\mathbf w)}\mid\mathbf v_i=a]\\
&=\frac1d\sum_{(a,b,c)\in E}x_ay_bz_c.\tag{9}
\end{align*}
Next, directly from the definition of$y$ and Claim24,
\[\mathbb E_{\mathbf w}[(-1)^{S_i(\mathbf w)}]=\frac1n\sum_{a\in[n]}y_a.\tag{10}\]
Similarly,
\[\mathbb E_{\mathbf w}[(-1)^{X_i(\mathbf w)}]=\frac1n\sum_{a\in[n]}x_a.\tag{11}\]
Equation(8) follows from Equations(9)–(11) and the fact that$\sum_{a\in[n]}z_a=1$. The desired result then follows Definition12.
\end{proof}
\subsection*{5.3 Decoding $C$}
We follow the work of Jernoimo et al. who gave a near-linear time list- and unique-decoding algorithm for Ta-Shma's code via an efficient weak regularity lemma, and prove that their algorithm also applies to our code as well. In the language of$\tau$-sampling distributions, they prove:\footnote{Jernoimo et al. prove their result under stronger requirements, however one can verify that the mixing requirement suffices.}
\paragraph{Theorem35 ([20]).} There exists a constant$c_{\mathrm{JST}}$ such that the following holds for any integers$d,t,k,n$ and any$\tau,\varepsilon_0,\varepsilon>0$. Let$C_0:\mathbb F_2^k\to\mathbb F_2^n$ be a code with bias at most$\varepsilon_0$ which is uniquely decodable to within distance$(1-\varepsilon_0)/4$ in time$T_0=T_0(n,\varepsilon_0)$. Let$\mathbf W\sim[n]^t$ be a homogeneous$\tau$-sampling distribution, let$C=\dsum_{\mathbf W}(C_0)$ be the corresponding direct sum lifting, and assume that the bias of$C$ is at most$\varepsilon$. Let$\beta$ be such that
\[\beta\geq\max\left\{\sqrt\varepsilon,\sqrt{c_{\mathrm{JST}}\cdot t^3\tau},2\left(\frac12+2\varepsilon_0\right)^t\right\}.\]
Then, there exists a randomized algorithm, which given$\tilde y\in\mathbb F_2^{|\mathbf W|}$, recovers the list$C\cap B(\tilde y,1/2-\beta)$ with probability at least$1-\frac1\varepsilon\cdot2^{-\Omega(\varepsilon_0^2n)}$ in time$\widetilde O(c_{\beta,t,\varepsilon_0}\cdot(|\mathbf W|+T_0))$, for$c_{\beta,t,\varepsilon_0}=(6/\varepsilon_0)^{2^{O(t^3/\beta^2)}}$. Moreover, if we are able to set$\beta=(1+\varepsilon)/4$, we can uniquely decode$C$ to within distance$(1-\varepsilon)/4$ with probability at least$1-2^{-\Omega(\varepsilon_0^2n)}$.
Plugging-in our code$C$ (using$\mathbf W_{H,t}$ and$C_0$ as described in Section5.1), we get our main result [Theorem36, stated above].
\begin{proof}[Proof of Theorem36]
By Lemma34, our parity sampler$\mathbf W$ which we use for$C$ is$\tau$-sampling for
\[\tau=\frac\lambda4=O\left(\frac{\log d}{\sqrt d}\right)=\widetilde O(2^{-2\sqrt{\log(1/\varepsilon)}}),\]
where we used the fact that our randomized construction of$H$ gives us$\lambda=O((1/\sqrt d)\log d)$ and the bound on$d$ from Equation(7). All that is left is to show how the two items follow from Theorem35. For the list decoding result, we take$\beta$ to be as small as possible. For us,
\[\left(\frac12+2\varepsilon_0\right)^t\leq2^{-\frac12\sqrt{\log(1/\varepsilon)}},\]
and$\sqrt{c_{\mathrm{JST}}\cdot t^3\tau}=\widetilde O(2^{-\sqrt{\log(1/\varepsilon)}})$. We can thus conclude that$\beta\leq2^{-\frac12\sqrt{\log(1/\varepsilon)}}$. (Again, we are assuming$\varepsilon$ is smaller than some small constant.) For the running time, note that
\[T_0=\exp(\operatorname{poly}(1/\varepsilon_0))n=\exp\left(2^{O(\sqrt{\log(1/\varepsilon)})}\right)\cdot k,\]
and$c_{\beta,t,\varepsilon_0}$ is thus triply-exponential in$\sqrt{\log(1/\varepsilon)}$, and overall$\widetilde O(c_{\beta,t,\varepsilon_0}\cdot(|\mathbf W|+T_0))=\widetilde O(c_{\beta,t,\varepsilon_0}\cdot k)$. For the unique decoding result, we take$\beta=(1+\varepsilon)/4$, and the running time becomes doubly-exponential in$\sqrt{\log(1/\varepsilon)}$.
\end{proof}

\section*{A Splittability and $\tau$-Sampling}
Our decoding result follows from the work of Jernoimo, Srivastava, and Tulsiani[20]. In[20], for the result of Theorem35 to hold, they require$\mathbf W$ to be$\tau$-splittable. $\tau$-splittability is stronger than, and in fact implies,$\tau$-sampling. We will not define splittability here, since for their result to hold, only a ``Splittable Mixing Lemma'' is required. The mixing property used in[20] goes as follows.
\paragraph{Definition56 (splittable mixing for $\pm1$ cut functions).} Given positive integers$n,t$, and$s<t-1$, we denote by$\mathcal F_s$ the set of all functions$f:[n]^t\to\{1,-1\}$ of the form
\[f(x_1,\ldots,x_t)=b\cdot\chi_{A_1}(x_1)\cdots\chi_{A_s}(x_s)\cdot\chi_B(x_{s+1},\ldots,x_t),\tag{27}\]
where$b\in\{1,-1\}$, $A_1,\ldots,A_s\subseteq[n]$, $B\subseteq[n]^{t-s}$, and where$\chi_S(x)=-1$ if$x\in S$ and1 otherwise.
For a distribution$\mathbf W\sim[n]^t$, we denote by$\mathbf W_{[i,j]}$ the distribution over$[n]^{j-i+1}$ that is obtained by sampling$\mathbf x\sim\mathbf W$ and outputting$\mathbf x_{[i,j]}$. For simplicity, we abbreviate$\mathbf W_i=\mathbf W_{[i,i]}$. Given$s\in[t-1]$, let$\nu_s$ be the probability measure over$[n]^t$ for which$\nu_s(x)=\prod_{i\in[s]}\Pr[\mathbf W_i=x_i]\cdot\Pr[\mathbf W_{[s+1,t]}=x_{[s+1,t]}]$. We say$\mathbf W\sim[n]^t$ satisfies splittable mixing with error$\tau$, if for every$s\in[t-1]$ and every$f,f'\in\mathcal F_s$ it holds that
\[\left|\mathbb E_{\mathbf x\sim\nu_s}[f(\mathbf x)f'(\mathbf x)]-\mathbb E_{\mathbf x\sim\nu_{s-1}}[f(\mathbf x)f'(\mathbf x)]\right|\leq\tau.\tag{28}\]
In[20], they show that the above property follows from$\tau$-splittablility. Here we show that the above property also follows from our weaker notion of$\tau$-sampling. For the sake of being compatible with[20], we will use a slightly different definition of$\tau$-sampling than what was given in Definition32.
\paragraph{Definition57.} We say that a distribution$\mathbf W$ over$[n]^t$ is$\tau$-sampling if for any$i\in[t-1]$, $S\subseteq[n]$, and$X\subseteq[n]^{t-i}$,
\[\Cov_{\mathbf w\sim\mathbf W}[1[\mathbf w_i\in S],1[(\mathbf w_{i+1},\ldots,\mathbf w_t)\in X]]\leq\tau.\]
For decoding lifted sum codes, both definitions are essentially the same, since we can always work with$\mathbf W_{\mathrm{rev}}$, wherein we sample$\mathbf w=(\mathbf w_1,\ldots,\mathbf w_t)\sim\mathbf W$ and output$(\mathbf w_t,\ldots,\mathbf w_1)$. Moreover, our$\mathbf W$ satisfies$\mathbf W=\mathbf W_{\mathrm{rev}}$.
\paragraph{Claim58.} Let$\mathbf W\subseteq[n]^t$ be homogeneous and$\tau$-sampling. Then,$\mathbf W$ satisfies splittable mixing with error$4\tau$.
\begin{proof}
Let$s\in[t-1]$, $b,b'\in\{1,-1\}$ and sets$A_1,\ldots,A_s,A'_1,\ldots,A'_s\subseteq[n]$ and$B,B'\subseteq[n]^{t-s}$ that correspond to the functions$f$ and$f'$. For$i\in[s]$ denote$C_i=A_i\triangle A'_i$ and$D=B\triangle B'$. We then have
\[f(x)f'(x)=\sigma\cdot\chi_{C_1}(x_1)\cdots\chi_{C_s}(x_s)\cdot\chi_D(x_{s+1},\ldots,x_t),\]
for$\sigma=bb'\in\{1,-1\}$. The subtraction on the left hand side of Equation(28) now reads
\begin{multline*}
\sum_{x\in[n]^{s-1}}\prod_{i\in[s-1]}\Pr[\mathbf W_i=x_i]\cdot\sigma\cdot\chi_{C_1}(x_1)\cdots\chi_{C_{s-1}}(x_{s-1})\cdot\\
\sum_{y\in[n],z\in[n]^{t-s}}\chi_{C_s}(y)\cdot\chi_D(z)\\
\cdot(\Pr[\mathbf W_s=y]\Pr[\mathbf W_{[s+1,t]}=z]-\Pr[\mathbf W_{[s,t]}=y\circ z]).\tag{29}
\end{multline*}
The second line of the above expression can be written as
\begin{multline*}
\mathbb E_{\mathbf y\sim\mathbf W_s}[(-1)^{1[\mathbf y\in C_s]}]\cdot\mathbb E_{\mathbf z\sim\mathbf W_{[s+1,t]}}[(-1)^{1[\mathbf z\in D]}]\\
-\mathbb E_{(\mathbf y,\mathbf z)\sim(\mathbf W_s,\mathbf W_{[s+1,t]})}[(-1)^{1[\mathbf y\in C_s]}(-1)^{1[\mathbf z\in D]}],\tag{30}
\end{multline*}
which is simply the covariance between the random variables$(-1)^{1[\mathbf y\in C_s]}$ and$(-1)^{1[\mathbf z\in D]}$ where$\mathbf y$ and$\mathbf z$ are drawn appropriately. As$\mathbf W$ is$\tau$-sampling, we know that
\[|\Cov[1[\mathbf y\in C_s],1[\mathbf z\in D]]|\leq\tau,\]
and so Equation(30), in absolute value, amounts to
\[|\Cov[(-1)^{1[\mathbf y\in C_s]},(-1)^{1[\mathbf z\in D]}]|\leq4\tau,\]
which readily follows from the fact that$(-1)^b=1-2b$ for$b\in\{0,1\}$, as we argued in Lemma34. Taking absolute value, by the triangle inequality, Equation(29) can be bounded by
\[4\tau\cdot\sum_{x\in[n]^{s-1}}\prod_{i\in[s-1]}\Pr[\mathbf W_i=x_i].\]
As$\mathbf W$ is homogeneous,$\Pr[\mathbf W_i=x_i]=1/n$, and we are done.
\end{proof}

\begin{thebibliography}{99}
\bibitem[1]{ref1} Noga Alon. Explicit expanders of every degree and size. Combinatorica, pages 1–17, 2021.
\bibitem[2]{ref2} Noga Alon, Jehoshua Bruck, Joseph Naor, Moni Naor, and Ron M. Roth. Construction of asymptotically good low-rate error-correcting codes through pseudo-random graphs. Informa- tion Theory, IEEE Transactions on, 38(2):509–516, 1992.
\bibitem[3]{ref3} Noga Alon, Oded Goldreich, Johan Håstad, and René Peralta. Simple constructions of almost k-wise independent random variables. Random Structures \& Algorithms, 3(3):289–304, 1992. New Near-Linear Time Decodable Codes Closer to the GV Bound
\bibitem[4]{ref4} Avraham Ben-Aroya and Amnon Ta-Shma. A combinatorial construction of almost-Ramanujan graphs using the zig-zag product. SIAM Journal on Computing, 40(2):267–290, 2011.
\bibitem[5]{ref5} Avraham Ben-Aroya and Amnon Ta-Shma. Constructing small-bias sets from algebraic- geometric codes. Theory of Computing, 9(5):253–272, 2013.
\bibitem[6]{ref6} Yonatan Bilu and Shlomo Hoory. On codes from hypergraphs. European Journal of Combin- atorics, 25(3):339–354, 2004.
\bibitem[7]{ref7} Yonatan Bilu and Nathan Linial. Lifts, discrepancy and nearly optimal spectral gap. Combin- atorica, 26(5):495–519, 2006.
\bibitem[8]{ref8} Andrej Bogdanov. A different way to improve the bias via expanders. Topics in (and out) the theory of computing, Lecture, 2012.
\bibitem[9]{ref9} Emma Cohen, Dhruv Mubayi, Peter Ralli, and Prasad Tetali. Inverse expander mixing for hypergraphs. The Electronic Journal of Combinatorics, 23(2):P2–20, 2016.
\bibitem[10]{ref10} David Conlon, Jonathan Tidor, and Yufei Zhao. Hypergraph expanders of all uniformities from Cayley graphs. Proceedings of the London Mathematical Society, 121(5):1311–1336, 2020.
\bibitem[11]{ref11} Joel Friedman and Avi Wigderson. On the second eigenvalue of hypergraphs. Combinatorica, 15(1):43–65, 1995.
\bibitem[12]{ref12} Edgar N. Gilbert. A comparison of signalling alphabets. The Bell system technical journal, 31(3):504–522, 1952.
\bibitem[13]{ref13} Oded Goldreich. On constructing expanders for any number of vertices, October 2019. Available at \url{https://www.wisdom.weizmann.ac.il/~oded/R1/ex4all.pdf}.
\bibitem[14]{ref14} Konstantin Golubev and Ori Parzanchevski. Spectrum and combinatorics of two-dimensional Ramanujan complexes. Israel Journal of Mathematics, 230(2):583–612, 2019.
\bibitem[15]{ref15} Venkatesan Guruswami and Piotr Indyk. Linear-time encodable/decodable codes with near- optimal rate. IEEE Transactions on Information Theory, 51(10):3393–3400, 2005.
\bibitem[16]{ref16} Venkatesan Guruswami, Atri Rudra, and Madhu Sudan. Essential Coding Theory, 2015. URL: \url{http://www.cse.buffalo.edu/faculty/atri/courses/coding-theory/book}.
\bibitem[17]{ref17} Brett Hemenway, Noga Ron-Zewi, and Mary Wootters. Local list recovery of high-rate tensor codes and applications. SIAM Journal on Computing, pages FOCS17–157, 2019.
\bibitem[18]{ref18} Wassily Hoeffding. Probability inequalities for sums of bounded random variables. Journal of the American statistical association, 58(301):13–30, 1963.
\bibitem[19]{ref19} Fernando Granha Jeronimo, Dylan Quintana, Shashank Srivastava, and Madhur Tulsiani. Unique decoding of explicit \(\varepsilon\)-balanced codes near the Gilbert–Varshamov bound. In Proceed- ings of the 61st Annual Symposium on Foundations of Computer Science (FOCS 2020), pages 434–445. IEEE, 2020.
\bibitem[20]{ref20} Fernando Granha Jeronimo, Shashank Srivastava, and Madhur Tulsiani. Near-linear time decoding of Ta-Shma’s codes via splittable regularity. In Proceedings of the 53rdth Annual Symposium on Theory of Computing (STOC 2021), pages 1527–1536. ACM, 2021.
\bibitem[21]{ref21} Swastik Kopparty, Nicolas Resch, Noga Ron-Zewi, Shubhangi Saraf, and Shashwat Silas. On list recovery of high-rate tensor codes. IEEE Transactions on Information Theory, 67(1):296–316, 2020.
\bibitem[22]{ref22} Swastik Kopparty, Noga Ron-Zewi, Shubhangi Saraf, and Mary Wootters. Improved decoding of folded Reed–Solomon and multiplicity codes. In Proceedings of the 59th Annual Symposium on Foundations of Computer Science (FOCS 2018), pages 212–223. IEEE, 2018.
\bibitem[23]{ref23} John Lenz and Dhruv Mubayi. Eigenvalues and linear quasirandom hypergraphs. In Forum of Mathematics, Sigma, volume 3. Cambridge University Press, 2015.
\bibitem[24]{ref24} Ray Li and Mary Wootters. Improved list-decodability of random linear binary codes. In APPROX–RANDOM, volume 116 of LIPIcs, pages 50:1–50:19. Schloss Dagstuhl - Leibniz- Zentrum für Informatik, 2018.
\bibitem[25]{ref25} Eyal Lubetzky and Yuval Peres. Cutoff on all Ramanujan graphs. Geometric and Functional Analysis, 26(4):1190–1216, 2016.
\bibitem[26]{ref26} Alexander Lubotzky, Ralph Phillips, and Peter Sarnak. Ramanujan graphs. Combinatorica, 8(3):261–277, 1988. 10:35
\bibitem[27]{ref27} Alexander Lubotzky, Beth Samuels, and Uzi Vishne. Explicit constructions of Ramanujan complexes of type Ad. European Journal of Combinatorics, 26(6):965–993, 2005.
\bibitem[28]{ref28} Grigorii Aleksandrovich Margulis. Explicit group-theoretical constructions of combinatorial schemes and their application to the design of expanders and concentrators. Problemy peredachi informatsii, 24(1):51–60, 1988.
\bibitem[29]{ref29} Jack Murtagh, Omer Reingold, Aaron Sidford, and Salil Vadhan. Deterministic approximation of random walks in small space. Theory of Computing, 17(1):1–35, 2021.
\bibitem[30]{ref30} Joseph Naor and Moni Naor. Small-bias probability spaces: Efficient constructions and applications. SIAM Journal on Computing, 22(4):838–856, 1993.
\bibitem[31]{ref31} Ori Parzanchevski. Mixing in high-dimensional expanders. Combinatorics, Probability and Computing, 26(5):746–761, 2017.
\bibitem[32]{ref32} Ori Parzanchevski, Ron Rosenthal, and Ran J. Tessler. Isoperimetric inequalities in simplicial complexes. Combinatorica, 36(2):195–227, 2016.
\bibitem[33]{ref33} Michael Sipser and Daniel A. Spielman. Expander codes. IEEE transactions on Information Theory, 42(6):1710–1722, 1996.
\bibitem[34]{ref34} Amnon Ta-Shma. Explicit, almost optimal, \(\varepsilon\)-balanced codes. In Proceedings of the 49th Annual Symposium on Theory of Computing (STOC 2017), pages 238–251. ACM, 2017.
\bibitem[35]{ref35} R. Tanner. A recursive approach to low complexity codes. IEEE Transactions on information theory, 27(5):533–547, 1981.
\bibitem[36]{ref36} Rom Rubenovich Varshamov. Estimate of the number of signals in error correcting codes. Docklady Akad. Nauk, SSSR, 117:739–741, 1957.
\bibitem[37]{ref37} Gillés Zémor. On expander codes. IEEE Transactions on Information Theory, 47(2):835–837, 2001.
\end{thebibliography}
\end{document}
